\chapter{Spécifier une fonctionnalité}\label{ch:spec}

\begin{objectifs}[Objectifs --- étape 2, \texttt{forge-spec}]
\begin{itemize}
  \item rédiger une spécification complète, règle par règle, avec la source de chaque règle ;
  \item écrire des scénarios Gherkin : nominal, limite, erreur ;
  \item distinguer ce que le script contrôle (le statut) de ce que vous devez relire ;
  \item obtenir une validation humaine, nommée et datée.
\end{itemize}
\end{objectifs}

\begin{situation}[Sur le terrain : \sika{}]
Première fonctionnalité du lot 1 : \q{Cotiser par Mobile Money}. Le développeur est tenté de commencer par
l'écran. La spec va d'abord forcer trois questions : quel montant est accepté ? quand une cotisation
compte-t-elle vraiment ? que se passe-t-il si un membre cotise deux fois ?
\end{situation}

\section{Le concept : une règle, une source, un exemple}
\begin{concept}[Spécification]
Le document qui décrit \emph{ce que la fonctionnalité doit faire}, sans dire comment. Chaque règle métier
y figure avec sa \textbf{source} (texte réglementaire, règlement intérieur, atelier) : une règle sans
source est une opinion.
\end{concept}
\begin{concept}[Gherkin]
Un langage d'exemples en français (\emph{Étant donné, Quand, Alors}). Un fichier \cmd{.feature}
contient une \cmd{Règle} par règle métier, et pour chacune au moins un cas nominal, un cas limite et un cas d'erreur.
Chaque scénario sert à la fois de critère d'acceptation et de test.
\end{concept}

\section{Pas à pas avec \texttt{forge-spec}}
\begin{enumerate}
  \item Le skill lit le glossaire, l'impact map et le gabarit, puis numérote la spec : \cmd{SPEC-001}.
  \item Il copie le gabarit en \cmd{docs/02-specs/SPEC-001.md}, statut \emph{Brouillon}.
  \item Il remplit section par section, une question à la fois : problème, récit, règles métier (chacune avec sa source), données (les données personnelles sont marquées), interfaces, hors périmètre.
  \item Il écrit \cmd{features/SPEC-001.feature}.
  \item Il passe la spec \emph{En revue} ; la validation, métier puis technique, est humaine.
\end{enumerate}
Le skill n'invente jamais une règle : quand elle manque, il l'inscrit en \emph{Questions ouvertes} et vous la demande.

\section{La spécification de \sika{} (extrait)}
L'extrait est condensé pour la lecture. Dans le fichier réel, l'en-tête est un tableau Markdown (\cmd{| Statut | Brouillon |}) : c'est cette ligne, et elle seule, que lit le garde-fou.

\begin{code}
\begin{Verbatim}
# SPEC-001 — Cotiser par Mobile Money

Statut : En revue · Contexte : Cotisation · Objectif : O1
Priorité : Must

## 2. Récit utilisateur
En tant que membre d'un groupe, je veux cotiser depuis mon
téléphone afin que ma cotisation soit enregistrée sans
contestation.

## 3. Règles métier
R1  Le montant est celui fixé pour le groupe.
    Source : règlement intérieur, art. 7
R2  Une cotisation ne compte qu'une fois confirmée par
    l'opérateur.
    Source : atelier Event Storming, point chaud H2
R3  Un membre cotise une seule fois par tour.
    Source : règlement intérieur, art. 8

## 5. Données
amountXof  entier, obligatoire, non personnelle
phone      texte, obligatoire, DONNÉE PERSONNELLE (APDP)

## 9. Questions ouvertes
Q1  Après combien de minutes sans réponse de l'opérateur une
    cotisation est-elle déclarée échouée ?
    Responsable : trésorière générale
\end{Verbatim}
\end{code}

\begin{piege}[Piège : valider avec une question ouverte bloquante]
Tant que Q1 n'a pas de réponse, R2 est incomplète : le cas limite \q{l'opérateur ne répond pas} ne peut
pas s'écrire. Le script, lui, ne le verra pas : il ne lit que le statut. C'est à vous de relire avant
de proposer \emph{Validée}. Ici, la trésorière générale répond \q{15 minutes} ; la question est close, et
le scénario limite peut être écrit.
\end{piege}

\subsection*{Les scénarios Gherkin}
\begin{code}
\begin{Verbatim}
# language: fr
Fonctionnalité: Cotiser par Mobile Money
  En tant que membre d'un groupe
  Je veux cotiser depuis mon téléphone
  Afin que ma cotisation soit enregistrée sans contestation

  Règle: R1 — le montant est celui fixé pour le groupe
    Scénario: Cotisation au montant du groupe
      Étant donné un groupe dont la cotisation est de 5000 FCFA
      Et un membre de ce groupe avec un tour ouvert
      Quand le membre cotise 5000 FCFA
      Alors la cotisation est initiée auprès de l'opérateur
    Scénario: Montant différent de celui du groupe
      Quand le membre cotise 3000 FCFA
      Alors l'action est refusée avec le motif "Montant incorrect"

  Règle: R2 — une cotisation ne compte qu'une fois confirmée
    Scénario: Confirmation par l'opérateur
      Étant donné une cotisation initiée de 5000 FCFA
      Quand l'opérateur confirme la cotisation
      Alors le solde du pot augmente de 5000 FCFA
      Et le membre reçoit un SMS de confirmation
    Scénario: Opérateur sans réponse après 15 minutes
      Étant donné une cotisation initiée depuis 15 minutes
      Quand le délai expire
      Alors la cotisation est échouée et ne compte pas dans le pot

  Règle: R3 — une seule cotisation par membre et par tour
    Scénario: Seconde cotisation dans le même tour
      Étant donné un membre ayant déjà cotisé pour ce tour
      Quand le membre cotise à nouveau
      Alors l'action est refusée avec le motif "Déjà cotisé pour ce tour"
\end{Verbatim}
\end{code}
Remarquez la lecture : le cas limite (15 minutes), le cas nominal et le cas d'erreur vivent côte à côte, sous la règle qu'ils éprouvent.
Remarquez aussi le vocabulaire : \emph{cotisation}, \emph{pot}, \emph{tour} --- pas \q{paiement}, \q{cagnotte}, \q{cycle}.

\section{Le garde-fou : le statut, et rien que le statut}
\begin{code}
\begin{Verbatim}
$ scripts/check-spec-validated.sh docs/02-specs/SPEC-001.md
KO: docs/02-specs/SPEC-001.md a le statut « Brouillon » (attendu : Validée)

$ scripts/check-spec-validated.sh docs/02-specs/SPEC-001.md
OK: docs/02-specs/SPEC-001.md est Validée
\end{Verbatim}
\end{code}
Le second résultat n'apparaît qu'après qu'un humain a écrit \textbf{Validée} (avec nom et date) dans le tableau d'en-tête.

\begin{retenir}[Avant de proposer « Validée », relisez cette liste]
\begin{itemize}
  \item les sections 1 à 5 sont renseignées ;
  \item aucun marqueur \cmd{[...]} ne subsiste ;
  \item aucune question ouverte bloquante ;
  \item le fichier \cmd{.feature} couvre chaque règle.
\end{itemize}
\end{retenir}

\begin{vousjouez}[À vous de jouer]
Choisissez la première fonctionnalité de votre projet. Écrivez deux règles métier avec leur source. Si l'une n'a pas de source,
écrivez-la en question ouverte et nommez qui peut y répondre. Rédigez trois scénarios (nominal, limite, erreur) en Gherkin français.
\end{vousjouez}

\begin{retenir}[À retenir]
\begin{itemize}
  \item Pas de code sans spec \emph{Validée}, et la validation est humaine.
  \item Règle $\Rightarrow$ source ; règle sans source $\Rightarrow$ question ouverte.
  \item Le script contrôle le statut ; la complétude se relit.
\end{itemize}
\end{retenir}
\source{skills/forge-spec/SKILL.md, templates/docs/02-specs}

\begin{acquis}
  \item Qu'écrit le skill quand une règle métier manque ?
  \item Que contrôle exactement \cmd{check-spec-validated.sh} ?
  \item Combien de scénarios minimum par règle, et de quels types ?
\end{acquis}
